% Generated draft: requires mathematical and source review.
\documentclass[11pt]{article}
\usepackage[margin=25mm]{geometry}
\usepackage{fontspec}
\setmainfont{Latin Modern Roman}
\usepackage{amsmath,amssymb,mathtools,unicode-math}
\setmathfont{Latin Modern Math}
\usepackage{enumitem,needspace,xurl,hyperref}
\setlength{\parindent}{0pt}
\setlength{\parskip}{5pt}
\setlength{\emergencystretch}{3em}
\newcommand{\N}{\mathbb{N}}
\newcommand{\Z}{\mathbb{Z}}
\newcommand{\Q}{\mathbb{Q}}
\newcommand{\R}{\mathbb{R}}
\newcommand{\C}{\mathbb{C}}
\newcommand{\F}{\mathbb{F}}
\newcommand{\A}{\mathbb{A}}
\newcommand{\PP}{\mathbb{P}}
\newcommand{\E}{\mathbb{E}}
\newcommand{\eps}{\varepsilon}
\newcommand{\dd}{\,\mathrm{d}}
\begin{document}
\section*{UnsolvedMath ID: 30003557}
\textbf{Source problem code:} catalogue entry
\subsection{Definitions and mathematical statement}
Let \(P=\R^2\setminus\{(0,0)\}\) denote the punctured plane, and let \(r\) tend to infinity through the positive integers. The supplied record asks whether
\[
\forall\varepsilon\in\R_{>0}\quad \exists\,G\text{ in }P\quad \exists\,L\in\R:\qquad
\lim_{r\to\infty}\tau(G_r)^{1/r}=L
\quad\text{and}\quad
1<L<1+\varepsilon.
\]
Here the graph \(G\) may depend on \(\varepsilon\).

This is a schematic statement: the supplied complete record does not define the graphs \(G_r\), the invariant \(\tau\), or the admissible class of graphs \(G\) in \(P\). Consequently, the displayed limit cannot yet be interpreted unambiguously. In particular, the record does not specify finiteness, connectedness, edge signs or weights, embedding requirements, or a covering construction for \(G_r\).

For terminology, a spanning tree of a finite graph \(H\) is a connected acyclic subgraph containing every vertex of \(H\). Its ordinary spanning-tree count is
\[
\tau_{\mathrm{st}}(H)=\#\{T\subseteq H:T\text{ is a spanning tree of }H\}.
\]
The supplied title suggests a relationship with spanning-tree counts, but the record does not establish that its symbol \(\tau\) means \(\tau_{\mathrm{st}}\); that identification is left unresolved.

The opening word ``Equivalently'' refers to an omitted preceding formulation. The supplied curation identifies a relationship with Lehmer's question, but neither the precise preceding question nor the hypotheses establishing the claimed equivalence are supplied. These qualifications must be resolved before this draft constitutes a fully specified mathematical problem.
\subsection{Short English statement}
The proposed question asks whether a graph in the punctured plane can have an associated exponential growth rate strictly between \(1\) and \(1+\varepsilon\), for every \(\varepsilon>0\). The graph family \(G_r\), the meaning of \(\tau\), and the allowed graphs are unspecified. The supplied record labels the problem open and asserts an equivalence with Lehmer's question; these are supplied claims requiring review.
\subsection{Sources}
\begin{itemize}
\item Supplied local UnsolvedMath registry record, permanent ID \(30003557\), title \emph{Spanning-Tree Growth Rates Near One}, exact source problem code \texttt{catalogue entry}, in the Oberwolfach Reports Open Problems collection. The record supplies proposed year 2017 and names no proposer.
\item Supplied report citation: \emph{Low-dimensional Topology and Number Theory} (2018), Oberwolfach Reports, pp.~2363--2426, DOI \texttt{10.4171/owr/2017/38}. Supplied source and literature link: \url{https://doi.org/10.4171/owr/2017/38}. The background describes Oberwolfach Reports as documenting mathematical workshops and their open-problem discussions. No problem-specific page or named contributor is supplied.
\item Supplied literature reference: \emph{Lehmer question and graph complexity}, \url{https://arxiv.org/abs/1106.1304}. Authors and a precise supporting passage are not supplied in the record.
\item Literature assessment dated exactly \texttt{2026-08-21}: the supplied curation labels the problem open, asserts that the graph reformulation is equivalent to Lehmer's question, and states that known special cases do not establish the asserted arbitrarily small gap above one. This is supplied dated curation derived from the supplied OWR curation, not an independently verified literature assessment. Its status and cited claims require independent verification.
\item The registry supplies statement status \texttt{corrected\_verified}, with an explanation that mathematical content was preserved while notation, typography, or omitted definitions were repaired. This is a supplied editorial label; the present record still leaves essential notation undefined, and no source was independently consulted for this draft.
\end{itemize}
\end{document}
